% ============================================================================
%  MadLyX shared LaTeX preamble
%
%  Every fix here is collected from "שימוש נכון בליך" (The MadLyX) by Kali,
%  with the page it comes from noted. Paste the whole thing into
%  Document > Settings > LaTeX Preamble, or keep only the parts you want.
%
%  The bundled templates already include this. This standalone copy is for
%  documents you started before installing MadLyX.
% ============================================================================


% ---------------------------------------------------------------------------
%  Muted colours                                                      (p.45-46)
%
%  LyX's stock colours are fluorescent and hard to read in print. These are the
%  same names, toned down, so existing coloured text just looks calmer.
%  Affects the PDF only; for the LyX editor window see \set_color in
%  preferences, which the installer sets on LyX 2.4.
% ---------------------------------------------------------------------------
\usepackage{xcolor}
\definecolor{blue}{RGB}{0, 0, 140}
\definecolor{red}{RGB}{170, 0, 0}
\definecolor{green}{RGB}{0, 110, 0}
\definecolor{magenta}{RGB}{150, 0, 150}
\definecolor{cyan}{RGB}{0, 120, 140}


% ---------------------------------------------------------------------------
%  Hebrew justification
%
%  Hebrew has essentially no hyphenation support in LaTeX, so TeX cannot break
%  long words to fill a line and you get "Overfull \hbox" warnings and text
%  spilling into the margin. Letting TeX stretch inter-word spacing further
%  fixes almost all of it.
%
%  Not from the guide - added by MadLyXInstaller.
% ---------------------------------------------------------------------------
\sloppy
\emergencystretch=3em


% ---------------------------------------------------------------------------
%  Blank first page                                                      (p.20)
%
%  Some documents render an empty first page before the content. Uncomment if
%  that happens to you. (Thanks to Amit Tsarfati, via the guide.)
% ---------------------------------------------------------------------------
% \setcounter{page}{0}
% \usepackage{atbegshi}
% \AtBeginDocument{\AtBeginShipoutNext{\AtBeginShipoutDiscard}}


% ---------------------------------------------------------------------------
%  Correct the \Longrightarrow tail                                   (p.71-72)
%
%  The stock \Longrightarrow renders with a malformed tail at some sizes.
%  This rebuilds it from a plain arrowhead and a rule.
% ---------------------------------------------------------------------------
\DeclareFontFamily{OT1}{cmrx}{}
\DeclareFontShape{OT1}{cmrx}{m}{n}{<->cmr10}{}
\let\saveLongrightarrow\Longrightarrow
\makeatletter
\renewcommand*{\Longrightarrow}{%
  \mathrel{\rlap{\fontfamily{cmrx}\fontencoding{OT1}\selectfont=}%
  \hphantom{\saveLongrightarrow}%
  \llap{$\m@th\Rightarrow$}}}
\makeatother

\let\saveLongleftarrow\Longleftarrow
\makeatletter
\renewcommand*{\Longleftarrow}{%
  \mathrel{\rlap{$\m@th\Leftarrow$}%
  \hphantom{\saveLongleftarrow}%
  \llap{\fontfamily{cmrx}\fontencoding{OT1}\selectfont =}}}
\makeatother


% ---------------------------------------------------------------------------
%  Disjoint union                                                         (p.72)
%
%  Gives \cupdot and \bigcupdot. To use them comfortably in LyX, define a math
%  macro (Insert > Math > Macro, or Alt+I H O) whose TeX side is \cupdot and
%  whose LyX side is \dot{\cup}.
% ---------------------------------------------------------------------------
\makeatletter
\def\moverlay{\mathpalette\mov@rlay}
\def\mov@rlay#1#2{\leavevmode\vtop{%
    \baselineskip\z@skip \lineskiplimit-\maxdimen
    \ialign{\hfil$\m@th#1##$\hfil\cr#2\crcr}}}
\newcommand{\charfusion}[3][\mathord]{
  #1{\ifx#1\mathop\vphantom{#2}\fi
     \mathpalette\mov@rlay{#2\cr#3}
  }
  \ifx#1\mathop\expandafter\displaylimits\fi}
\makeatother
\newcommand{\cupdot}{\charfusion[\mathbin]{\cup}{\cdot}}
\newcommand{\bigcupdot}{\charfusion[\mathop]{\bigcup}{\cdot}}


% ---------------------------------------------------------------------------
%  Indicator function                                                     (p.73)
%
%  The guide answers this one repeatedly: the indicator is \mathds{1}.
% ---------------------------------------------------------------------------
\usepackage{dsfont}
\newcommand{\indicator}{\mathds{1}}


% ---------------------------------------------------------------------------
%  Code listings                                                       (p.63-64)
%
%  Styling for Insert > Program Listing. mathescape lets you write maths inside
%  a listing by wrapping it in $...$. (Thanks to Ziv Rosenwasser, via the guide.)
% ---------------------------------------------------------------------------
\usepackage{listings}
\definecolor{codegreen}{rgb}{0, 0.6, 0}
\definecolor{codegray}{rgb}{0.5, 0.5, 0.5}
\definecolor{codepurple}{rgb}{0.58, 0, 0.82}
\definecolor{backcolour}{rgb}{0.97, 0.98, 0.99}
\definecolor{codeblue}{rgb}{0, 0, 0.3}
\lstdefinestyle{mystyle}{
  backgroundcolor=\color{backcolour},
  commentstyle=\color{codegreen},
  keywordstyle=\color{codeblue},
  numberstyle=\tiny\color{codegray},
  stringstyle=\color{codepurple},
  basicstyle=\ttfamily\footnotesize,
  breakatwhitespace=false,
  breaklines=true,
  captionpos=b,
  keepspaces=true,
  numbersep=5pt,
  showspaces=false,
  showstringspaces=false,
  showtabs=false,
  tabsize=2,
  mathescape=true
}
\lstset{style=mystyle}


% ---------------------------------------------------------------------------
%  Two-column sections                                                    (p.73)
%
%  For a document that is one column in places and two in others, leave
%  Document > Settings > Text Layout on one column and wrap the two-column part
%  in \begin{multicols}{2} ... \end{multicols} using Ctrl+L (TeX Code).
% ---------------------------------------------------------------------------
\usepackage{multicol}


% ---------------------------------------------------------------------------
%  Clickable cross-references                                          (p.53-54)
%
%  Turns cross-references into links and generates a PDF bookmark tree.
%
%  The guide (p.53) warns that hyperref errors in the Hebrew Article class.
%  That no longer reproduces: tested against babel[hebrew] + culmus with
%  current MiKTeX, hyperref compiles cleanly, alone and alongside the packages
%  below. Enabled by default; comment it out if your document disagrees.
% ---------------------------------------------------------------------------
\usepackage{hyperref}
\hypersetup{colorlinks=true, linkcolor=blue, urlcolor=blue, bookmarks=true}


% ===========================================================================
%  OPTIONAL EXTRAS
%
%  None of these are in the guide. Each was compiled against a Hebrew document
%  before being listed here, because RTL handling conflicts with a fair number
%  of LaTeX packages and guessing is unreliable. The installer already puts
%  them on your machine, so uncommenting a line is all that is needed.
%
%  Known NOT to work with Hebrew: algorithmicx (algorithm + algpseudocode).
%  Use algorithm2e below instead.
% ===========================================================================


% --- Physics and maths notation --------------------------------------------
%  \dv{f}{x} derivative, \pdv partial, \ket{\psi} \bra{\phi}, \abs{x}, \norm{v}
%  Probably the single biggest quality-of-life win for a physics student.
% \usepackage{physics}

%  Units that space and format themselves: \SI{9.81}{\meter\per\second\squared}
% \usepackage{siunitx}

%  Dirac notation, if you prefer it to the physics package: \Braket{\phi|\psi}
% \usepackage{braket}

%  \diff{f}{x} as an alternative derivative syntax.
% \usepackage{esdiff}

%  \mathscr{L} script letters, and \mathbbm{1} for the indicator.
% \usepackage{mathrsfs}
% \usepackage{bbm}


% --- Presentation ----------------------------------------------------------
%  Coloured, breakable boxes - the nicest way to make theorems stand out.
%  \begin{tcolorbox}[colback=blue!5,title=משפט] ... \end{tcolorbox}
% \usepackage[most]{tcolorbox}

%  Cross-references that name themselves: \cref{thm:main} prints "Theorem 3".
%  Must be loaded after hyperref.
% \usepackage{cleveref}

%  Subtle micro-typography. Worth having in Hebrew, where the absence of
%  hyphenation makes justification harder.
% \usepackage{microtype}

%  Control over list spacing: \begin{itemize}[noitemsep]
% \usepackage{enumitem}

%  Tables that do not look like a spreadsheet, and ones that fit a width.
% \usepackage{booktabs}
% \usepackage{tabularx}

%  Figure captions and side-by-side subfigures.
% \usepackage{caption}
% \usepackage{subcaption}
% \usepackage{wrapfig}


% --- Diagrams and algorithms -----------------------------------------------
%  Draw anything; plot functions and data.
% \usepackage{tikz}
% \usepackage{pgfplots}
% \pgfplotsset{compat=1.18}

%  Pseudocode. Use this rather than algorithmicx, which breaks under Hebrew.
% \usepackage[ruled,vlined]{algorithm2e}
